\section{Contexto del sistema}

\subsection{PuellaGame como organización y como sistema de información}

PuellaGame es un centro de entretenimiento familiar que atraviesa un periodo de
alta afluencia. La llegada de numerosos visitantes durante las vacaciones de
verano vuelve especialmente visible una necesidad que ya existía: conservar de
manera ordenada la información con la que opera el negocio. Aunque la solución
definitiva será una base de datos, su diseño y construcción requieren tiempo. El
prototipo planteado en esta práctica ocupa, por tanto, una posición intermedia:
permite comenzar a registrar datos en archivos CSV sin perder de vista que estos
serán utilizados posteriormente para poblar la base de datos.

Esta circunstancia define el sentido del sistema. No se trata solamente de
guardar texto en archivos, sino de establecer una primera representación del
negocio. Cada dato capturado debe pertenecer a una entidad reconocible, conservar
un significado estable y poder recuperarse mediante una llave. De este modo, el
prototipo funciona al mismo tiempo como herramienta operativa inmediata y como
puente hacia una solución de persistencia más completa.

\begin{figure}[htbp]
  \centering
  \setlength{\fboxsep}{7pt}
  \begin{tabular}{c@{\hspace{0.8em}$\longrightarrow$\hspace{0.8em}}c@{\hspace{0.8em}$\longrightarrow$\hspace{0.8em}}c}
    \fbox{\begin{minipage}{0.22\textwidth}\centering
      \textbf{Mundo real}\\[2pt]
      clientes, sucursales\\
      y premios
    \end{minipage}}
    &
    \fbox{\begin{minipage}{0.25\textwidth}\centering
      \textbf{Prototipo}\\[2pt]
      captura, valida y\\
      permite operar
    \end{minipage}}
    &
    \fbox{\begin{minipage}{0.22\textwidth}\centering
      \textbf{Persistencia}\\[2pt]
      archivos CSV listos\\
      para su migración
    \end{minipage}}
  \end{tabular}
  \caption{El prototipo como enlace entre la operación cotidiana y la futura
  base de datos.}
  \label{fig:puellagame-puente}
\end{figure}

\subsection{Actores y relaciones}

El administrador de PuellaGame es el actor que origina la necesidad. Su interés
no está en los detalles internos de los archivos, sino en disponer cuanto antes
de información persistente, consultable y suficientemente organizada para no
comenzar desde cero cuando se construya la base de datos. Desde su perspectiva,
el sistema debe ofrecer una imagen confiable de las sucursales, los premios y
los clientes registrados.

La interacción diaria recae en el personal autorizado que opera el menú del
prototipo. Esta persona traduce hechos del entorno en registros: da de alta una
sucursal, incorpora un premio, captura los datos de un cliente, localiza un
registro por su llave y corrige o elimina información cuando la realidad cambia.
El operador es, por ello, el punto de contacto entre la organización y el
sistema; la claridad de los mensajes y la validación de las entradas resultan
esenciales para que esa traducción sea correcta.

El cliente participa de forma indirecta. Sus datos forman parte del dominio,
pero no se establece que utilice personalmente la aplicación. Proporciona la
información necesaria durante la atención en el centro y el personal la incorpora
al sistema. Las sucursales y los premios tampoco son usuarios: son objetos del
negocio cuya existencia debe representarse de forma consistente. Esta distinción
evita confundir a quien realiza una acción con aquello sobre lo que la acción se
realiza.

\begin{figure}[htbp]
  \centering
  \setlength{\tabcolsep}{3pt}
  \renewcommand{\arraystretch}{1.25}
  \begin{tabular}{>{\centering\arraybackslash}p{0.20\textwidth}
                  >{\centering\arraybackslash}p{0.04\textwidth}
                  >{\centering\arraybackslash}p{0.25\textwidth}
                  >{\centering\arraybackslash}p{0.04\textwidth}
                  >{\centering\arraybackslash}p{0.20\textwidth}}
    \fbox{\parbox{0.17\textwidth}{\centering\textbf{Administrador}\\define la necesidad\\y consulta resultados}}
    & $\leftrightarrow$
    & \fbox{\parbox{0.22\textwidth}{\centering\textbf{Personal autorizado}\\opera el menú y recibe\\la respuesta del sistema}}
    & $\leftrightarrow$
    & \fbox{\parbox{0.17\textwidth}{\centering\textbf{Prototipo}\\valida, procesa y\\persiste en CSV}}
    \\[1.2em]
    \multicolumn{5}{c}{\small\textbf{Entidades representadas:} sucursal, premio y cliente}
  \end{tabular}
  \caption{Mapa de actores y elementos que rodean al prototipo. Las flechas
  representan intercambio de información, no acceso directo a los archivos.}
  \label{fig:mapa-contexto}
\end{figure}

\subsection{Flujo de trabajo observado}

El ciclo comienza cuando ocurre un hecho relevante para PuellaGame: se abre o
modifica una sucursal, se incorpora un premio o se registra un cliente. El
operador selecciona en el menú la entidad y la operación correspondiente, e
introduce la información solicitada. Antes de alterar la persistencia, el
prototipo comprueba que los campos esperados como numéricos realmente contengan
números y que la operación pueda ejecutarse. Si encuentra un problema, conserva
el estado previo e informa la causa mediante un mensaje comprensible.

Cuando la entrada es válida, el sistema localiza el archivo asociado con la
entidad y realiza el cambio. Una consulta sigue el recorrido inverso: a partir
de la llave proporcionada, el prototipo busca el registro y devuelve todos sus
datos al operador. La edición sustituye la representación anterior por la nueva,
mientras que la eliminación retira el registro identificado. En todos los casos,
el archivo CSV constituye la memoria del prototipo: permite que la información
permanezca disponible después de cerrar el programa.

\begin{figure}[htbp]
  \centering
  \small
  \renewcommand{\arraystretch}{1.45}
  \begin{tabularx}{0.96\textwidth}{@{}>{\bfseries}p{0.16\textwidth}X>{\centering\arraybackslash}p{0.21\textwidth}@{}}
    \toprule
    Momento & Transformación de la información & Resultado visible \\
    \midrule
    Entrada & El operador elige una entidad y expresa una alta, consulta,
    modificación o baja mediante el menú. & Datos capturados \\
    Control & El prototipo interpreta la solicitud, valida los valores y maneja
    cualquier condición excepcional. & Confirmación o aviso \\
    Persistencia & La representación de la entidad se lee o se actualiza en el
    archivo CSV correspondiente. & Registro recuperable \\
    \bottomrule
  \end{tabularx}
  \caption{Recorrido de la información dentro del contexto del sistema.}
  \label{fig:recorrido-informacion}
\end{figure}

\subsection{Frontera y alcance}

La frontera del sistema separa las decisiones humanas y los acontecimientos del
negocio de las responsabilidades propias del software. Fuera de ella quedan la
decisión de abrir una sucursal, la definición comercial de un premio, la atención
al cliente y la determinación de qué registros deben cambiar. Dentro de ella se
encuentran la recepción de entradas, su validación, la ejecución de las cuatro
operaciones de mantenimiento, la búsqueda por llave, el manejo de errores y la
lectura y escritura de los archivos CSV.

La futura base de datos también queda fuera del alcance inmediato. El prototipo
no administra concurrencia entre múltiples usuarios, no define todavía un motor
de base de datos y no sustituye el modelado posterior. Sin embargo, sí debe
preparar ese paso: separar las tres clases de entidad, conservar llaves estables
y evitar que los archivos contengan valores imposibles de interpretar. La
calidad de esta primera captura determinará cuánto trabajo de limpieza será
necesario durante la migración.

Visto en conjunto, PuellaGame no necesita simplemente tres archivos, sino un
flujo controlado que mantenga correspondencia entre la actividad real y su
representación digital. El administrador obtiene continuidad para la operación,
el personal dispone de un medio uniforme para trabajar y la organización genera
un acervo inicial que podrá evolucionar hacia una base de datos. Esa relación
entre personas, actividades, entidades y persistencia constituye el contexto
completo del sistema propuesto.
